% !TeX root = ../main.tex
% Section 3: No-Arbitrage and the Quantum Pricing Formula
% Batch #4.3 — Post M1+M2 axis, notation hygiene applied

\section{No-Arbitrage and the Quantum Pricing Formula}\label{sec:no-arbitrage}

The Weyl-Heisenberg algebra of Section~\ref{sec:weyl-heisenberg} provides the
algebraic foundation for a rigorous formulation of the hedging problem under price
impact. We now show that minimizing quadratic hedging error in this non-commutative
setting is equivalent to an orthogonal projection in the GNS Hilbert space, yielding
a closed-form quantum pricing formula that reduces to the classical Black--Scholes
rule in the commutative limit $\hbar_{\mathrm{econ}} \to 0$.

\begin{assumption}[Market state in the hedgeable subspace]\label{ass:phi-in-K}
The market risk state $\Phi_t$ lies in the hedgeable subspace $\mathcal{K}_t$:
$\Phi_t \in \mathcal{K}_t$. This regularity condition ensures that the
market state itself can be expressed as a limit of tradeable portfolio positions.
It is formally justified and verified for the single-qubit parametrisation
in Section~\ref{sec:duality-hedging} (Assumption~\ref{ass:phi-in-K-detail}).
\end{assumption}

\subsection{The Hedging Problem with Price Impact}

In frictionless BSM markets, any contingent claim can be perfectly replicated by a
self-financing portfolio. This relies on the assumption that hedging trades do not
affect prices. In real markets, every hedge trade alters the underlying price,
creating a feedback loop that we formalize through the non-commutative algebra of
market observables.

Let $\mathcal{A}$ be the $C^*$-algebra generated by the log-price operator $\hat{X}$
and the order-flow momentum operator $\hat{P}$, satisfying the Weyl form of the
canonical commutation relation \eqref{eq:weyl-ccr}. The market state is described by
a faithful normal state $\omega$ on $\mathcal{A}$, which we identify with the GNS
cyclic vector $|\Omega_\omega\rangle$ in the representation
$\pi_\omega : \mathcal{A} \to \mathcal{B}(\mathcal{H}_\omega)$.

\begin{definition}[Hedgeable subspace, restated]\label{def:hedgeable-subspace-t}
Let $\mathcal{K} \subset \mathcal{H}_\omega$ be the closed linear subspace spanned
by the discounted asset prices available for hedging:
\begin{equation*}
    \mathcal{K} = \overline{\operatorname{span}}
        \{|S_{t_i}\rangle : t_i \in [0, T]\},
\end{equation*}
where $|S_{t_i}\rangle$ is the GNS vector corresponding to the discounted asset
price at time $t_i$. We call $\mathcal{K}$ the \emph{hedgeable subspace}.
\end{definition}

\begin{lemma}[Non-triviality of the hedgeable subspace]\label{lem:K-nontrivial}
Under the Stone--von Neumann representation
$\mathcal{H}_\omega \simeq L^2(\mathbb{R},\mathbb{C})$ of
Proposition~\ref{prop:stone-vn}, and assuming the discounted asset price process
$(S_t)_{t \in [0,T]}$ is non-degenerate (i.e., $S_t$ is not almost-surely constant),
the hedgeable subspace $\mathcal{K}$ satisfies
$\{0\} \subsetneq \mathcal{K} \subsetneq \mathcal{H}_\omega$.
\end{lemma}

\begin{proof}
\emph{Non-emptiness} ($\mathcal{K} \neq \{0\}$): Each $|S_{t_i}\rangle$ is the
GNS image of the non-zero observable $\hat{S}_{t_i} \in \mathcal{A}$; faithfulness
of $\omega$ gives $\omega(\hat{S}_{t_i}^* \hat{S}_{t_i}) > 0$, so
$\||S_{t_i}\rangle\|^2 = \omega(\hat{S}_{t_i}^*\hat{S}_{t_i}) > 0$.

\emph{Properness} ($\mathcal{K} \neq \mathcal{H}_\omega$): The wavefunction of
$|S_{t_i}\rangle$ in the log-price representation is localised near the asset
level $S_{t_i}$ (it belongs to a compact subset of the spectrum of $\hat{X}$).
The payoff $g(x)=(e^x - K)_+$ of a call option grows without bound as
$x \to +\infty$, so $|g\rangle = \pi_\omega(\hat{G})|\Omega_\omega\rangle \notin
\mathcal{K}$ for generic $K > 0$. Hence $\mathcal{K}$ is a proper subspace,
confirming market incompleteness for $\hbar_{\mathrm{econ}} > 0$.
\end{proof}

The hedging problem is to find a portfolio $|h\rangle \in \mathcal{K}$ that
minimizes the quadratic hedging error:
\begin{equation}\label{eq:hedging-error}
    \min_{|h\rangle \in \mathcal{K}} \, \| |Y\rangle - |h\rangle \|_\omega^2,
\end{equation}
where $|Y\rangle = e^{-rT} |G(S_T)\rangle$ is the discounted payoff vector.

\begin{theorem}[Existence and uniqueness of the optimal hedge]\label{thm:existence-hedge}
The hedging problem \eqref{eq:hedging-error} has a unique solution
$|h^*\rangle \in \mathcal{K}$, given by the orthogonal projection of $|Y\rangle$
onto $\mathcal{K}$:
\begin{equation}\label{eq:optimal-hedge}
    |h^*\rangle = \Pi_{\mathcal{K}} |Y\rangle,
\end{equation}
where $\Pi_{\mathcal{K}} \in \mathcal{B}(\mathcal{H}_\omega)$ is the orthogonal
projection operator onto $\mathcal{K}$.
\end{theorem}

\begin{proof}
The functional $J(h) = \|Y - h\|_\omega^2$ is strictly convex and coercive on the
closed subspace $\mathcal{K}$. By the Hilbert space projection theorem, there exists
a unique minimizer $|h^*\rangle \in \mathcal{K}$ satisfying the orthogonality
condition
\begin{equation*}
    \langle Y - h^* | k \rangle_\omega = 0 \quad \forall \, |k\rangle \in \mathcal{K}.
\end{equation*}
This condition is equivalent to $|h^*\rangle = \Pi_{\mathcal{K}} |Y\rangle$.
\end{proof}

\begin{corollary}[Minimum hedging error]\label{cor:min-hedging-error}
The minimum achievable hedging error is
\begin{equation}\label{eq:min-hedging-error}
    \varepsilon_{\min}^2 = \|Y\|_\omega^2 - \|\Pi_{\mathcal{K}} Y\|_\omega^2.
\end{equation}
In the classical frictionless limit $\hbar_{\mathrm{econ}} \to 0$, $\mathcal{K}$
becomes the full space $\mathcal{H}_\omega$ and $\varepsilon_{\min} = 0$, recovering
BSM perfect replication.
\end{corollary}

\subsection{Hedging as Quantum Measurement}

Equation \eqref{eq:hedging-error} is equivalent to a quantum measurement: the
hedging portfolio projects the payoff onto the subspace of hedgeable claims
\citep{vonneumann1932mathematical}.

\begin{definition}[Hedging dual vector]\label{def:hedging-dual}
Let $|h^*\rangle = \Pi_{\mathcal{K}} |Y\rangle$ be the optimal hedge vector. The
\emph{hedging dual vector} $\langle h^* |$ is the unique element of the dual space
$\mathcal{H}_\omega^*$ obtained via the Riesz representation theorem: for any
$|\psi\rangle \in \mathcal{H}_\omega$,
$\langle h^* | \psi \rangle = \langle h^* | \psi \rangle_\omega$.
\end{definition}

\begin{proposition}[Hedging as a quantum measurement]\label{prop:hedging-measurement}
The optimal hedge $|h^*\rangle = \Pi_{\mathcal{K}} |Y\rangle$ is the
orthogonal projection of the payoff-state vector onto the hedgeable subspace:
\begin{equation*}
    |h^*\rangle = \Pi_{\mathcal{K}}\,|Y\rangle
    \;=\; \Pi_{\mathcal{K}}\,\pi_\omega(\hat{Y})|\Omega_\omega\rangle,
\end{equation*}
where $\Pi_{\mathcal{K}} : \mathcal{H}_\omega \to \mathcal{K}$ is the
orthogonal projection operator onto the hedgeable subspace and
$|Y\rangle = \pi_\omega(\hat{Y})|\Omega_\omega\rangle$ is the GNS
representative of the discounted payoff. At the operator level, the optimal hedge
satisfies $\hat{h}^* = \mathbb{E}_\omega[\hat{Y}\,|\,\mathcal{A}_{\mathcal{K}}]$
where $\mathbb{E}_\omega[\cdot|\mathcal{A}_{\mathcal{K}}]$ is Takesaki's
conditional expectation onto the sub-algebra $\mathcal{A}_{\mathcal{K}}$
associated with $\mathcal{K}$
\citep{takesaki2003theory}; at the GNS vector level this is equivalent to
$|h^*\rangle = \Pi_{\mathcal{K}}|Y\rangle$ as above.
\end{proposition}

\begin{proof}[Proof (deferred to Theorem~\ref{thm:hedging-measurement} in \S\ref{sec:duality-hedging})]
The result follows as a corollary of the Hedging--Measurement Duality
(Theorem~\ref{thm:hedging-measurement} in Section~\ref{sec:duality-hedging}),
which is proved independently of this proposition.
That theorem establishes
$\pi_\omega(\hat{h}^*)|\Omega_\omega\rangle
 = \Pi_{\mathcal{K}}\,\pi_\omega(\hat{Y})|\Omega_\omega\rangle$
in the GNS representation, from which taking the inner product with
$|\Omega_\omega\rangle$ (and using
$\omega(A) = \langle\Omega_\omega|\pi_\omega(A)|\Omega_\omega\rangle$) yields
the trace formula for $|h^*\rangle$.
\end{proof}

\subsection{The Main Pricing Formula}

\begin{theorem}[Quantum pricing formula]\label{thm:pricing}
Let $G(S_T)$ be the payoff of a European option maturing at $T$, let $r \geq 0$ be
the risk-free rate, and let
$|\Phi_t\rangle \in \mathcal{H}_\omega \otimes \mathcal{H}_{\mathrm{anc}}$ be the
purified market risk state at time $t \leq T$, evolving under the Lindblad master
equation \eqref{eq:lindblad}. Then the arbitrage-free price of the option at time
$t$ is
\begin{equation}\label{eq:pricing-master}
    P_t = \langle h^* \,|\, \Phi_t \rangle \, e^{-r(T-t)},
\end{equation}
where:
\begin{enumerate}
    \item[(i)] $|Y\rangle = e^{-rT} |G(S_T)\rangle$ is the discounted payoff vector;
    \item[(ii)] $|h^*\rangle = \Pi_{\mathcal{K}} |Y\rangle$ is the optimal hedge
        vector, the orthogonal projection of $|Y\rangle$ onto $\mathcal{K}$;
    \item[(iii)] $\langle h^* |$ is the Riesz dual vector of $|h^*\rangle$;
    \item[(iv)] $e^{-r(T-t)}$ is the deterministic discount factor.
\end{enumerate}
Corollary~\ref{cor:state-hedge-pricing} specializes this formula to the real Hilbert
space $\mathcal{H}_\omega \simeq L^2(\mathbb{R}, \mathbb{R})$, making realness and
positivity manifest and providing the operational form used for empirical
calibration.
\end{theorem}

\begin{proof}
The proof proceeds in five steps.

\paragraph{Step 1: GNS representation of the payoff.}
By the GNS construction, the discounted payoff $Y = e^{-rT} G(S_T)$ is represented
as $|Y\rangle = \pi_\omega(e^{-rT} \hat{G}) |\Omega_\omega\rangle$.

\paragraph{Step 2: Optimal hedge as orthogonal projection.}
By Theorem~\ref{thm:existence-hedge}, $|h^*\rangle = \Pi_{\mathcal{K}} |Y\rangle$.

\paragraph{Step 3: Duality pairing and the market risk state.}
The market risk state $|\Phi_t\rangle$ purifies $\rho_t$:
$\rho_t = \operatorname{Tr}_{\mathrm{anc}}(|\Phi_t\rangle \langle \Phi_t|)$. The
value of the optimal hedge at time $t$ is
$V_t^{\mathrm{hedge}} = \operatorname{Tr}(\rho_t \hat{h}^*)
= \langle \Phi_t | \hat{h}^* \otimes \mathbf{1}_{\mathrm{anc}} | \Phi_t \rangle
= \langle h^* | \Phi_t \rangle$.

\paragraph{Step 4: Discounting.}
$P_t = e^{-r(T-t)} V_t^{\mathrm{hedge}}$.

\paragraph{Step 5: No-arbitrage consistency.}
The market pricing functional
$\omega_t(A) = \langle \Phi_t | \pi_\omega(A) | \Phi_t \rangle$ is positive,
normalized, and faithful. By Theorem~\ref{thm:fundamental}, this is equivalent to
the absence of arbitrage.
\end{proof}

\begin{remark}[Relation to the BSM formula]\label{rem:pricing-bsm-limit}
In the classical limit $\hbar_{\mathrm{econ}} \to 0$, $[\hat{X}, \hat{P}] = 0$, the
GNS Hilbert space reduces to $L^2(\mathbb{P})$, $\mathcal{K} = \mathcal{H}_\omega$,
$\Pi_{\mathcal{K}} = \mathbf{1}$, and $|h^*\rangle = |Y\rangle$. The pricing formula
collapses to
$P_t^{\mathrm{BSM}} = e^{-r(T-t)} \mathbb{E}^{\mathbb{Q}}[G(S_T) | \mathcal{F}_t]$,
recovering the classical BSM price.
\end{remark}

\subsection{The No-Arbitrage Condition in the Quantum Framework}

The fundamental theorem of quantum asset pricing extends the classical results of
\citet{harrison1979martingales} and \citet{delbaen1994general} to the non-commutative
setting: it states that a faithful normal state on the observable algebra is
equivalent to the absence of quantum arbitrage.

\begin{definition}[Quantum arbitrage]\label{def:quantum-arbitrage}
A \emph{quantum arbitrage} is a non-zero observable $A \in \mathcal{A}$ such that
$A \geq 0$ and $\omega(A) = 0$ for some market state $\omega$. In other words, an
arbitrage is a non-negative financial position that has zero cost but is not
identically zero.
\end{definition}

\begin{theorem}[Fundamental theorem of quantum asset pricing]\label{thm:fundamental}
Under Assumption~\ref{ass:phi-in-K} ($\Phi_t \in \mathcal{K}_t$), the following
are equivalent:
\begin{enumerate}
    \item[(i)] There exists a faithful normal state $\omega$ on $\mathcal{A}$.
    \item[(ii)] The market is free of quantum arbitrage.
    \item[(iii)] The GNS representation $\pi_\omega$ is faithful and the pricing
        functional $P_t = \langle h^* | \Phi_t \rangle e^{-r(T-t)}$ is strictly
        positive for all non-zero non-negative payoffs.
\end{enumerate}
\end{theorem}

\begin{proof}
We prove the equivalence in three steps.

\paragraph{(i) $\Rightarrow$ (ii).}
Suppose $\omega$ is faithful and $A \geq 0$ with $\omega(A) = 0$. By faithfulness,
$A = 0$, so no non-zero arbitrage exists.

\paragraph{(ii) $\Rightarrow$ (i).}
Assume the market is free of quantum arbitrage (condition (ii)).
By Online Supplement~\S{S.13} (the NFLVR--GNS equivalence theorem therein), the NFLVR
condition---of which no-quantum-arbitrage is the quantum analogue---implies
that the market state $\omega$ is a faithful normal state on $\mathcal{A}$.
Specifically, the argument in Online Supplement~\S{S.13} constructs $\omega$
via Gleason's theorem on the Type~I von Neumann algebra $\pi_\omega(\mathcal{A})''
= \mathcal{B}(\mathcal{H}_\omega)$: complete additivity of the valuation
functional on projections (ensured by no-arbitrage) yields a unique
trace-class density operator $\rho$ with $\omega = \mathrm{Tr}(\rho\,\cdot)$
and $\rho > 0$ (faithfulness). This establishes~(i).

\begin{remark}[Gleason's theorem and the two-dimensional case]
\label{rem:gleason-dim}
Gleason's theorem \citep{gleason1957measures} requires $\dim(\mathcal{H}) \geq 3$.
For the single-qubit parametrisation ($\dim = 2$) used in the empirical
Section~\ref{sec:empirical-calibration}, the Bloch-sphere representation
provides a direct alternative: every state on $\mathcal{B}(\mathbb{C}^2)$
has the form $\omega(A) = \tfrac{1}{2}\operatorname{Tr}((I + \mathbf{r}\cdot\boldsymbol{\sigma})A)$
with $\|\mathbf{r}\|\leq 1$, and faithfulness is equivalent to $\|\mathbf{r}\|<1$.
The absence of arbitrage in the two-dimensional case is verified directly:
$\omega(A) = 0$ for $A \geq 0$ implies $\operatorname{Tr}(\rho A) = 0$, which
with $\rho > 0$ forces $A = 0$. Theorem~\ref{thm:fundamental} therefore holds
for all $\dim(\mathcal{H}) \geq 2$, with the Gleason argument applying when
$\dim \geq 3$ and the explicit Bloch construction covering $\dim = 2$.
\end{remark}

\paragraph{(i) $\Leftrightarrow$ (iii).}
If $\omega$ is faithful and normal, the GNS representation
$\pi_\omega : \mathcal{A} \to \mathcal{B}(\mathcal{H}_\omega)$ is injective
(faithful). Strict positivity of $P_t$ follows because for any non-zero
non-negative payoff $\hat{G} \geq 0$, faithfulness gives
$\omega(\hat{G}) > 0$, and by Assumption~\ref{ass:phi-in-K}
($\Phi_t \in \mathcal{K}_t$) together with the positivity argument of
Remark~\ref{rem:state-hedge-econ},
$P_t = \langle h^* | \Phi_t \rangle e^{-r(T-t)}
= \langle Y_t | \Phi_t \rangle e^{-r(T-t)} > 0$
(since $\varepsilon_t \perp \mathcal{K}_t$ by construction and
$\Phi_t \in \mathcal{K}_t$ by Assumption~\ref{ass:phi-in-K},
so $\varepsilon_t \perp \Phi_t$).
Conversely, if $P_t > 0$ for all non-zero non-negative payoffs, then
$\omega(\hat{G}) > 0$ for all such $\hat{G}$, which is precisely
faithfulness of $\omega$.
\end{proof}

\begin{corollary}[No-arbitrage and the economic Planck constant]\label{cor:no-arb-hbar}
For $\hbar_{\mathrm{econ}} > 0$, the Weyl algebra $\mathcal{A}$ is simple:
it has no non-trivial closed two-sided ideals. Consequently, every non-zero
$*$-representation of $\mathcal{A}$ is faithful (injective). In particular,
the GNS representation $\pi_\omega$ for any non-zero state $\omega$ is
faithful, so the pricing formula is well-defined.

\begin{remark}
Simplicity of $\mathcal{A}$ implies that every \emph{representation} is
faithful, but not that every \emph{state} is faithful in the sense that
$\omega(A^*A) = 0 \Rightarrow A = 0$. A vector state
$\omega_\psi(A) = \langle\psi|A|\psi\rangle$ satisfies faithfulness only if
$\psi$ is a separating vector (i.e.\ $A\psi = 0 \Rightarrow A = 0$). For
the Schr\"{o}dinger representation on $L^2(\mathbb{R})$, any non-zero $\psi$
is separating for bounded operators, so vector states are faithful in
practice. The faithful normal state whose existence is guaranteed by
Theorem~\ref{thm:fundamental}(i) is a generic element of the state space
that satisfies this separating property by construction.
\end{remark}

In the classical limit $\hbar_{\mathrm{econ}} \to 0$, the algebra becomes
commutative and faithfulness of $\omega$ reduces to the classical condition
that the pricing measure be equivalent to the physical measure.
\end{corollary}

\begin{proof}
The simplicity of the Weyl $C^*$-algebra for $\hbar_{\mathrm{econ}} \neq 0$
is a standard result \citep{petz1990algebra}. Injectivity of every non-zero
representation follows directly from simplicity: if
$\ker(\pi) \neq \{0\}$, then $\ker(\pi)$ is a non-trivial closed two-sided
ideal, contradicting simplicity.
\end{proof}

\subsection{The Gibbs State and the KMS Condition}

In thermal equilibrium, quantum systems are described by Gibbs states. The
financial analogue is a market in which the pricing functional is determined by a
Hamiltonian $\hat{H}$ and an inverse temperature $\beta = 1/(k_B T)$, where $k_B$
is the Boltzmann constant and $T$ is the ``market temperature.''

\begin{definition}[Gibbs market state]\label{def:gibbs-state}
A \emph{Gibbs market state} at inverse temperature $\beta > 0$ is a state
$\omega_\beta$ on $\mathcal{A}$ defined by
\begin{equation*}
    \omega_\beta(A) = \frac{\operatorname{Tr}(e^{-\beta \hat{H}} A)}
        {\operatorname{Tr}(e^{-\beta \hat{H}})}, \quad A \in \mathcal{A},
\end{equation*}
where $\hat{H}$ is the market Hamiltonian \eqref{eq:hamiltonian}.
\end{definition}

The Gibbs state satisfies the Kubo--Martin--Schwinger (KMS) condition, which
characterizes thermal equilibrium in quantum statistical mechanics
\citep{kubo1957statistical,martin1959theory}. In the financial context, the KMS
condition implies that the market is in a state of ``no asymptotic arbitrage'' in
the sense that no trading strategy can generate risk-free profits over infinitely
long horizons.

\begin{theorem}[KMS condition and no asymptotic arbitrage]\label{thm:kms}
Let $\omega_\beta$ be a Gibbs market state at inverse temperature $\beta$. Then for
any $A, B \in \mathcal{A}$, there exists a function $F_{A,B}(z)$ analytic in the
strip $0 < \Im(z) < \beta$ and continuous on its closure such that
\begin{equation*}
    F_{A,B}(t) = \omega_\beta(A \alpha_t(B)), \qquad
    F_{A,B}(t + i\beta) = \omega_\beta(\alpha_t(B) A),
\end{equation*}
where $\alpha_t(A) = e^{it\hat{H}} A e^{-it\hat{H}}$ is the Heisenberg time
evolution. This condition is equivalent to the absence of asymptotic arbitrage.
\end{theorem}

\begin{proof}
The KMS condition is a standard result for Gibbs states: for
$\omega_\beta(A) = Z^{-1}\mathrm{Tr}(e^{-\beta\hat{H}}A)$, the function
$F_{A,B}(z) = Z^{-1}\mathrm{Tr}(e^{-\beta\hat{H}}A\,e^{iz\hat{H}}B\,e^{-iz\hat{H}})$
is analytic in the strip $\{0 < \Im(z) < \beta\}$ and satisfies the
boundary conditions by the cyclicity of the trace \citep{bratteli1981operator}.

For the equivalence with no asymptotic arbitrage: an asymptotic arbitrage
is a sequence of strategies $\{H_n\} \subset \mathcal{A}$ with
$\omega_\beta(H_n^*H_n) \to 0$ (vanishing cost) but
$\liminf_{n\to\infty} \omega_\beta((H_n - \varepsilon)_+) > 0$
for some $\varepsilon > 0$ (bounded profit). If such a sequence existed,
then for large $n$ we have $H_n \approx 0$ in the $L^2(\omega_\beta)$
sense, so $F_{H_n,B}(z) = \omega_\beta(H_n\,\alpha_z(B)) \to 0$
uniformly on compacts by the Cauchy estimate for bounded analytic functions.
But the boundary value $F_{H_n,B}(t+i\beta) = \omega_\beta(\alpha_t(B)H_n)
\to \omega_\beta(\alpha_t(B) \cdot 0) = 0$ as well. Since $\omega_\beta(H_n^*H_n)\to 0$
implies $H_n \to 0$ in the $L^2(\omega_\beta)$-norm (weak operator topology),
the continuous functional calculus is weakly continuous
\citep{bratteli1981operator}: $f(H_n)\to f(0)$ weakly for every continuous $f$.
Applying $f(x)=(x-\varepsilon)_+$ gives $(H_n-\varepsilon)_+ \to 0$ weakly,
hence $\omega_\beta((H_n-\varepsilon)_+) \to 0$, contradicting
$\liminf\,\omega_\beta((H_n-\varepsilon)_+) > 0$. Therefore no asymptotic
arbitrage can exist for any Gibbs state $\omega_\beta$.
\end{proof}

\begin{remark}[Economic interpretation of the KMS condition]\label{rem:kms-econ}
The KMS condition has a natural economic interpretation: it states that the market
is in a state of ``thermal equilibrium'' where the pricing functional is
time-translation invariant. In this equilibrium, the price of an option is
determined by the Hamiltonian $\hat{H}$ and the inverse temperature $\beta$, which
together encode the market's risk preferences and the cost of microstructural
friction. The KMS condition also provides a rigorous foundation for the use of the
Lindblad master equation \eqref{eq:lindblad} in modeling the evolution of the
market state: the Gibbs state is the stationary solution of the Lindblad dynamics
in the absence of external driving.
\end{remark}

\subsection{The Pricing Formula in the Gibbs State}

When the market is in a Gibbs state, the pricing formula takes a particularly
elegant form. The optimal hedge vector $|h^*\rangle$ can be expressed directly in
terms of the Hamiltonian and the inverse temperature.

\begin{proposition}[Gibbs state pricing]\label{prop:gibbs-pricing}
Let $\omega_\beta$ be a Gibbs market state at inverse temperature $\beta$,
defined by $\omega_\beta(A) = Z^{-1}\operatorname{Tr}(e^{-\beta\hat{H}}A)$
where $\hat{H}$ is the market Hamiltonian and $Z = \operatorname{Tr}(e^{-\beta\hat{H}})$.
Then the optimal hedge vector in the pricing formula
\eqref{eq:pricing-master} is
\begin{equation*}
    |h^*\rangle = \Pi_{\mathcal{K}} |Y\rangle,
\end{equation*}
where $|Y\rangle$ is the discounted payoff vector. This formulation
guarantees that the inner product $\langle h^* | \Phi_t \rangle$ is
manifestly real and non-negative for any non-negative payoff.
\end{proposition}

\begin{proof}
For a KMS state $\omega_\beta$ with respect to the time-evolution
automorphism group $\alpha_t(A) = e^{it\hat{H}}Ae^{-it\hat{H}}$,
the Tomita--Takesaki modular operator satisfies $\Delta = e^{-\beta\hat{H}}$
\citep{bratteli1981operator}, i.e., the modular operator coincides with
the Gibbs weight. In the GNS representation of $\omega_\beta$, the cyclic
vector $|\Omega_\beta\rangle$ is a separating vector for
$\pi_{\omega_\beta}(\mathcal{A})$, and the modular conjugation $J$ satisfies
$J\pi_\omega(A)J = \pi_\omega(A^*)^*$ on $\mathcal{H}_{\omega_\beta}$.
The optimal hedge $|h^*\rangle = \Pi_{\mathcal{K}}|Y\rangle$ follows from
Theorem~\ref{thm:hedging-measurement}; the Gibbs structure of $\omega_\beta$
ensures that the KMS analyticity condition holds (see also
Theorem~\ref{thm:kms}), which guarantees both the faithfulness of
$\omega_\beta$ (by Theorem~\ref{thm:fundamental}) and the strict positivity
of the pricing formula.

\begin{remark}
The equation $\pi_\omega(e^{-\beta\hat{H}})|\Omega_\omega\rangle = |h^*\rangle$
holds specifically when the payoff operator $\hat{G}$ commutes with the
modular flow, i.e., when $\hat{G}$ is in the centraliser of $\omega_\beta$.
For general payoffs, $|h^*\rangle$ is obtained by projection as above.
\end{remark}
\end{proof}

\subsection{Dynamics of the Pricing Formula}

The pricing formula \eqref{eq:pricing-master} evolves in time as the market risk
state $|\Phi_t\rangle$ evolves under the Lindblad master equation. This evolution
provides a dynamic hedging strategy that is consistent with the no-arbitrage
condition.

\begin{proposition}[Dynamics of the option price]\label{prop:price-dynamics}
The option price $P_t = \langle h^* | \Phi_t \rangle e^{-r(T-t)}$ evolves according
to
\begin{equation*}
    dP_t = r P_t \, dt + \langle h^* | d\Phi_t \rangle e^{-r(T-t)},
\end{equation*}
where $d\Phi_t$ is the stochastic differential of the market risk state under the
Lindblad dynamics \eqref{eq:lindblad}. The hedge ratio is given by the duality
pairing $\Delta_t = \langle h^* | \partial_S \Phi_t \rangle e^{-r(T-t)}$.
\end{proposition}

\begin{proof}
In this proposition $h^*$ is treated as the time-$t$ reference hedge, fixed at
the projection $h^* = \Pi_{\mathcal{K}_t}Y_t$ (i.e., recalibrated at each $t$ but
treated as constant for the infinitesimal increment). Differentiating
$P_t = \langle h^* | \Phi_t \rangle e^{-r(T-t)}$ with respect to $t$ yields the
discount term $rP_t\,dt$ plus the stochastic increment
$\langle h^* | d\Phi_t \rangle e^{-r(T-t)}$, where $d\Phi_t$ is governed by the
It\^o form of the Lindblad master equation. The hedge ratio
$\Delta_t = \langle h^* | \partial_S \Phi_t \rangle e^{-r(T-t)}$ follows from the
chain rule applied to the $S$-dependence of $\Phi_t$. The full derivation of the
It\^o--Lindblad stochastic differential and the delta formula is given in
Online Supplement~S.1.2.
\end{proof}
